% GENERATED FILE. Edit canonical lessons, then run npm run generate:books.
% canonical-source-hash: fnv1a64:984725d209d00296
% canonical-lessons: UR-C472-hazm-karna, UR-C472-hazma, UR-C472-badhazmi, UR-C472-ubkai, UR-C472-ghiza

\chapter{An Upset Stomach}
\label{ch:ur-an-upset-stomach}
\hlchaptermodality{472}

\begin{chapteropening}
\textbf{By the end of this chapter:} \emph{I can say that someone has an upset stomach.}

Complete the last lesson of chapter 472: I can say that someone has an upset stomach.
\end{chapteropening}

\section[\texorpdfstring{\ur{ہضم} \ur{کرنا} (hazm karnā)}{ہضم کرنا (hazm karnā)}]{\ur{ہضم} \ur{کرنا} (hazm karnā) — to digest}

\label{lesson:UR-C472-hazm-karna}



\begin{quote}
Before the new one: say the Urdu for healing, then the Urdu for medical.
\end{quote}

\subsection*{You'll want to know: \ur{ہضم} \ur{کرنا}}
\textbf{\ur{ہضم} \ur{کرنا}} — \emph{hazm karnā} — \textquotedblleft{}to digest\textquotedblright{}.

\begin{quote}
\ur{معدہ} \ur{کھانا} \ur{ہضم} \ur{کرتا} \ur{ہے۔}
\end{quote}

\begin{quote}
\emph{medā khānā hazm kartā hai.}
\end{quote}

\begin{quote}
The stomach digests food.
\end{quote}

Past, with \ur{نے} (he or she …): \ur{اس} \ur{نے} \ur{کھانا} \ur{ہضم} \ur{کیا}, \emph{us ne khānā hazm kiyā}.

\begin{cousinweb}
\textbf{\ur{ہضم} \ur{کرنا}} (\emph{hazm karnā}) is built from \textbf{\ur{ہضم}}, usually traced, through Persian, to Arabic \emph{haḍm}, \textquotedblleft{}digestion\textquotedblright{}, and \textbf{\ur{کرنا}}, \textquotedblleft{}to do\textquotedblright{}.
\end{cousinweb}

\begin{grammarlens}[title={What you've built}]
One more word for describing what you lived and felt.
\end{grammarlens}

\subsection*{Your turn}
\begin{itemize}
\raggedright
  \item \emph{Say it:} \emph{hazm karnā}
  \item \emph{Say it:} \emph{hazm karnā}, once more
  \item \emph{Say it:} \emph{medā khānā hazm kartā hai.} — The stomach digests food.
  \item \emph{Recall:} say the Urdu for healing, then the Urdu for medical, then say \emph{hazm karnā} again
\end{itemize}

\begin{tcolorbox}[breakable,colback=teal!4,colframe=teal!35!black,title={Before you move on}]
What does \ur{ہضم} \ur{کرنا} mean? (\textquotedblleft{}To digest\textquotedblright{}.) Say it once more.
\end{tcolorbox}
\section[\texorpdfstring{\ur{ہاضمہ} (hāzmā)}{ہاضمہ (hāzmā)}]{\ur{ہاضمہ} (hāzmā) — digestion}

\label{lesson:UR-C472-hazma}



\begin{quote}
Before the new one: say the Urdu for medical, then the Urdu for to digest.
\end{quote}

\subsection*{You'll want to know: \ur{ہاضمہ}}
\textbf{\ur{ہاضمہ}} — \emph{hāzmā} — \textquotedblleft{}digestion\textquotedblright{}.

\begin{quote}
\ur{دہی} \ur{ہاضمے} \ur{کے} \ur{لیے} \ur{اچھا} \ur{ہے۔}
\end{quote}

\begin{quote}
\emph{dahī hāzme ke liye achchhā hai.}
\end{quote}

\begin{quote}
Yoghurt is good for the digestion.
\end{quote}

Masculine.

\begin{cousinweb}
\textbf{\ur{ہاضمہ}} (\emph{hāzmā}) is usually traced, through Persian, to Arabic \emph{hāḍima}, \textquotedblleft{}digesting\textquotedblright{}, of the same family as \textbf{\ur{ہضم}}.
\end{cousinweb}

\begin{grammarlens}[title={What you've built}]
One more word for describing what you lived and felt.
\end{grammarlens}

\subsection*{Your turn}
\begin{itemize}
\raggedright
  \item \emph{Say it:} \emph{hāzmā}
  \item \emph{Say it:} \emph{hāzmā}, once more
  \item \emph{Say it:} \emph{dahī hāzme ke liye achchhā hai.} — Yoghurt is good for the digestion.
  \item \emph{Recall:} say the Urdu for medical, then the Urdu for to digest, then say \emph{hāzmā} again
\end{itemize}

\begin{tcolorbox}[breakable,colback=teal!4,colframe=teal!35!black,title={Before you move on}]
What does \ur{ہاضمہ} mean? (\textquotedblleft{}Digestion\textquotedblright{}.) Say it once more.
\end{tcolorbox}
\section[\texorpdfstring{\ur{بدہضمی} (badhazmī)}{بدہضمی (badhazmī)}]{\ur{بدہضمی} (badhazmī) — indigestion}

\label{lesson:UR-C472-badhazmi}



\begin{quote}
Before the new one: say the Urdu for to digest, then the Urdu for digestion.
\end{quote}

\subsection*{You'll want to know: \ur{بدہضمی}}
\textbf{\ur{بدہضمی}} — \emph{badhazmī} — \textquotedblleft{}indigestion\textquotedblright{}.

\begin{quote}
\ur{زیادہ} \ur{کھانے} \ur{سے} \ur{بدہضمی} \ur{ہو} \ur{گئی۔}
\end{quote}

\begin{quote}
\emph{zyādā khāne se badhazmī ho gaī.}
\end{quote}

\begin{quote}
Eating too much gave me indigestion.
\end{quote}

Feminine.

\begin{cousinweb}
\textbf{\ur{بدہضمی}} (\emph{badhazmī}) is usually traced to Persian \emph{badhazmī}, built from \emph{bad}, \textquotedblleft{}bad\textquotedblright{}, and \textbf{\ur{ہضم}}, \textquotedblleft{}digestion\textquotedblright{}, from Arabic \emph{haḍm}, with the ending \emph{-ī}.
\end{cousinweb}

\begin{grammarlens}[title={What you've built}]
One more word for describing what you lived and felt.
\end{grammarlens}

\subsection*{Your turn}
\begin{itemize}
\raggedright
  \item \emph{Say it:} \emph{badhazmī}
  \item \emph{Say it:} \emph{badhazmī}, once more
  \item \emph{Say it:} \emph{zyādā khāne se badhazmī ho gaī.} — Eating too much gave me indigestion.
  \item \emph{Recall:} say the Urdu for to digest, then the Urdu for digestion, then say \emph{badhazmī} again
\end{itemize}

\begin{tcolorbox}[breakable,colback=teal!4,colframe=teal!35!black,title={Before you move on}]
What does \ur{بدہضمی} mean? (\textquotedblleft{}Indigestion\textquotedblright{}.) Say it once more.
\end{tcolorbox}
\section[\texorpdfstring{\ur{ابکائی} (ubkāī)}{ابکائی (ubkāī)}]{\ur{ابکائی} (ubkāī) — retching; feeling sick}

\label{lesson:UR-C472-ubkai}



\begin{quote}
Before the new one: say the Urdu for digestion, then the Urdu for indigestion.
\end{quote}

\subsection*{You'll want to know: \ur{ابکائی}}
\textbf{\ur{ابکائی}} — \emph{ubkāī} — \textquotedblleft{}retching; feeling sick\textquotedblright{}.

\begin{quote}
\ur{بس} \ur{میں} \ur{اسے} \ur{ابکائی} \ur{آئی۔}
\end{quote}

\begin{quote}
\emph{bas meṅ use ubkāī āī.}
\end{quote}

\begin{quote}
She felt sick on the bus.
\end{quote}

Feminine.

\begin{cousinweb}
\textbf{\ur{ابکائی}} (\emph{ubkāī}) is inherited: it is the noun of the Hindi verb \emph{ubaknā}, \textquotedblleft{}to vomit\textquotedblright{}, which CDIAL lists under \emph{ubbakka}, \textquotedblleft{}vomited\textquotedblright{}, a form it reconstructs (CDIAL 2337).
\end{cousinweb}

\begin{grammarlens}[title={What you've built}]
One more word for describing what you lived and felt.
\end{grammarlens}

\subsection*{Your turn}
\begin{itemize}
\raggedright
  \item \emph{Say it:} \emph{ubkāī}
  \item \emph{Say it:} \emph{ubkāī}, once more
  \item \emph{Say it:} \emph{bas meṅ use ubkāī āī.} — She felt sick on the bus.
  \item \emph{Recall:} say the Urdu for digestion, then the Urdu for indigestion, then say \emph{ubkāī} again
\end{itemize}

\begin{tcolorbox}[breakable,colback=teal!4,colframe=teal!35!black,title={Before you move on}]
What does \ur{ابکائی} mean? (\textquotedblleft{}Retching; feeling sick\textquotedblright{}.) Say it once more.
\end{tcolorbox}
\section[\texorpdfstring{\ur{غذا} (ghizā)}{غذا (ghizā)}]{\ur{غذا} (ghizā) — food; nourishment}

\label{lesson:UR-C472-ghiza}



\begin{quote}
Before the new one: say the Urdu for indigestion, then the Urdu for retching.
\end{quote}

\subsection*{You'll want to know: \ur{غذا}}
\textbf{\ur{غذا}} — \emph{ghizā} — \textquotedblleft{}food; nourishment\textquotedblright{}.

\begin{quote}
\ur{مریض} \ur{کو} \ur{ہلکی} \ur{غذا} \ur{دیجیے۔}
\end{quote}

\begin{quote}
\emph{marīz ko halkī ghizā dījiye.}
\end{quote}

\begin{quote}
Give the patient light food.
\end{quote}

Feminine.

\begin{cousinweb}
\textbf{\ur{غذا}} (\emph{ghizā}) is usually traced, through Persian, to Arabic \emph{ghidhāʾ}, \textquotedblleft{}nourishment\textquotedblright{}.
\end{cousinweb}

\begin{grammarlens}[title={What you've built}]
One more word for describing what you lived and felt.
\end{grammarlens}

\subsection*{Your turn}
\begin{itemize}
\raggedright
  \item \emph{Say it:} \emph{ghizā}
  \item \emph{Say it:} \emph{ghizā}, once more
  \item \emph{Say it:} \emph{marīz ko halkī ghizā dījiye.} — Give the patient light food.
  \item \emph{Recall:} say the Urdu for indigestion, then the Urdu for retching, then say \emph{ghizā} again
\end{itemize}

\begin{tcolorbox}[breakable,colback=teal!4,colframe=teal!35!black,title={Before you move on}]
What does \ur{غذا} mean? (\textquotedblleft{}Food; nourishment\textquotedblright{}.) Say it once more.
\end{tcolorbox}
